\documentclass[11pt,a4paper]{article}
\usepackage[utf8]{inputenc}
\usepackage[T1]{fontenc}
\usepackage{lmodern}
\usepackage{textcomp}
\usepackage{amsmath,amssymb}
\usepackage{graphicx}
\usepackage{booktabs}
\usepackage{longtable}
\usepackage{array}
\usepackage{listings}
\usepackage{xcolor}
\usepackage[margin=1in]{geometry}
\usepackage{float}
\usepackage[hidelinks,breaklinks]{hyperref}
\usepackage{microtype}
\setlength{\parskip}{4pt}
\lstset{basicstyle=\ttfamily\scriptsize,breaklines=true,breakatwhitespace=false,
  postbreak=\mbox{\textcolor{gray}{$\hookrightarrow$}\space},frame=single,
  framerule=0.3pt,rulecolor=\color{gray!60},columns=fullflexible,keepspaces=true,
  showstringspaces=false,numbers=none,xleftmargin=4pt,xrightmargin=4pt,
  literate={\_}{{\_}}1}
\graphicspath{{figures/}}
\providecommand{\tightlist}{}
\title{Sterile-Neutrino Dark Matter and Big-Bang Nucleosynthesis Constraints}
\author{HEP-Genesis-Agent session record \\ \small Project: hep-testing-research}
\date{2026-10-15}
\begin{document}
\maketitle
\noindent\textit{\small This document is an AI-generated account of an interactive session with the HEP-Genesis agent. Its text was written by a language model strictly from the session record; figures and numbers come only from tool results obtained during the session. The supplementary material and the transcript, generated alongside, document exactly what ran. Verify independently before use.}
\vspace{1em}

\begin{abstract}
The session surveyed arXiv literature on sterile-neutrino dark matter, Big-Bang nucleosynthesis (BBN), and heavy neutral leptons (HNLs), separating the keV dark-matter regime from the MeV--GeV HNL regime. The directly relevant minimal BBN analysis, arXiv:2410.07343, parameterizes the HNL cosmology by mass and flavor-specific active mixing and includes temporary matter domination together with meson, photon, and neutrino injection. Its qualitative BBN exclusion is an intermediate-mixing regime, but numerical electron-, muon-, and tau-mixing contour coordinates were not recovered. For keV dark matter, the record supports only conditional viability: minimal non-resonant production supplying all dark matter is excluded by X-ray constraints in the cited self-interaction study, while resonant, mixed-component, self-interaction, and scalar-mediated scenarios can reopen parameter space under model-dependent assumptions. Quantitative examples include a 7.1-keV mixed model with astrophysical mixing near \(7\times10^{-10}\) and viable values below roughly \(10^{-9}\), and a refined resonant-production region with \(m_s\gtrsim20\,\mathrm{keV}\) and lepton-to-entropy ratio at least of order \(10^{-3}\). These results are conditional rather than universal mass--mixing bounds.
\end{abstract}

\section{Introduction}

The session began with a request to survey recent arXiv papers from 2024--2026 on sterile-neutrino dark matter and BBN constraints, and to identify the mass--mixing regions that remain viable. The investigation then expanded to compare keV sterile-neutrino dark matter with heavier HNLs, because the two regimes are constrained by different cosmological effects. In the keV regime, the relevant issues were relic abundance, radiative-decay X-rays, momentum distributions, free streaming, and small-scale structure. In the MeV--GeV regime, the focus was production, lifetime, temporary matter domination, decay products, and BBN energy injection.

The session also examined resonant Shi--Fuller production, including a 7.1-keV mixed cold-plus-warm benchmark and a refined analysis reporting viable regions above approximately \(20\,\mathrm{keV}\). Additional directions concerned non-standard neutrino self-interactions, scalar-mediated freeze-in/freeze-out, axion-like-particle (ALP) cascades, and HNL dark decays. The record does not contain a completed numerical scan combining all of these constraints, nor does it establish universal mass--mixing intervals independent of the production model.

\section{Methods}

The session used local arXiv abstract and metadata searches, direct arXiv searches, full-text downloads, and chunk-based reading. The execution dispatch was local. No input files or datasets were present at the beginning of the session. The focused searches included queries for sterile-neutrino dark matter with BBN constraints, keV sterile-neutrino dark matter with primordial nucleosynthesis, HNL BBN decays, X-ray and Lyman-\(\alpha\) constraints, non-standard neutrino self-interactions, and scalar freeze-in/freeze-out. The autonomous literature-survey route attempted coverage through 2026-10-14, but full-text extraction failed with the error \texttt{RuntimeError: 'Search' object has no attribute 'results'} for its selected papers.

Direct full-text downloads succeeded for arXiv:2306.16532, arXiv:2112.00758, arXiv:2204.08795, arXiv:2410.07343, arXiv:2507.12270, arXiv:2410.06970, and arXiv:2510.01907. The available record includes chunked full-text reading for arXiv:2306.16532, arXiv:2112.00758, arXiv:2204.08795, arXiv:2410.07343, and arXiv:1702.08430. The record does not include completed full-text extraction or figure-caption verification for all downloaded papers. No new numerical code, BBN network calculation, parameter scan, lifetime calculation, or abundance calculation was run during the session. The supplement associated with this report contains the recorded tool-call trail and available extracted text; it does not supply missing contour coordinates or unperformed calculations.

The evidence board classified most substantive findings as contested, despite assigning high confidence to their source-level summaries. Accordingly, the results below distinguish reported claims from independently reproduced calculations. The numerical values quoted from literature abstracts or extracted text retain the assumptions of those sources.

\section{Results}

\subsection{Literature coverage and regime separation}

The focused search identified arXiv:2410.07343 as the directly relevant minimal-model BBN analysis. The search also returned 2024--2025 studies of HNL dark decays, ALP cascades, and resonant keV sterile-neutrino production. No clearly relevant 2025--2026 paper combining standard keV sterile-neutrino dark matter with a dedicated BBN abundance analysis was returned by the recorded focused searches. This is a property of the recorded search results, not a proof that no such paper exists.

The literature separates into two regimes. The keV studies concern dark-matter production and its momentum distribution, while the HNL studies concern cosmological decay and energy injection. The 2017 review arXiv:1702.08430 identifies non-resonant, resonant, and decay production as distinct mechanisms and emphasizes that the momentum distribution is needed to assess structure formation. The session therefore did not treat a BBN-allowed HNL region as automatically relevant to keV dark matter.

\subsection{Minimal heavy-neutrino BBN constraints}

The analysis arXiv:2410.07343, ``BBN Constraint on Heavy Neutrino Production and Decay,'' uses the two parameters
\[
\left(m_4,\lvert U_{a4}\rvert^2\right),
\qquad a=e,\mu,\tau .
\]
According to the extracted abstract and the board finding, the calculation includes early-Universe production, temporary matter domination, and energy injection through charged mesons, photons, and light neutrinos. It presents BBN constraints as complementary to terrestrial HNL searches.

The qualitative exclusion structure reported in the session is an intermediate-mixing window. At sufficiently large mixing, the heavy neutrino decays before it significantly affects BBN. At sufficiently small mixing, freeze-in production is too weak to alter the light-element abundances. Very long-lived states are subject to a separate CMB constraint. The extracted text also reports that, above the pion-mass threshold, the lifetime contours turn downward because the total decay rate increases and the production mechanism changes rapidly from freeze-out to freeze-in.



Figure 1 (not in record) represents the published BBN-excluded heavy-neutrino mass--mixing plane of arXiv:2410.07343, including the temporary-matter-domination and decay-injection effects described above. It bears directly on the question of viable mass--mixing regions by showing that the BBN-affected domain is not simply a one-sided upper or lower mixing limit. However, the session did not recover numerical contour coordinates for the electron-, muon-, or tau-mixing cases. The figure therefore supports only the qualitative exclusion structure in this report, not a reproducible table of flavor-specific boundaries. This finding was marked contested by the research board because the contours were not independently extracted during the session.

\subsection{Thermally equilibrated sub-GeV HNLs and dark decays}

The 2025 paper arXiv:2507.12270 reports that thermally equilibrated sub-GeV HNLs require
\[
\tau_N\lesssim 0.02\,\mathrm{s}
\]
for successful BBN. Its abstract states that significant dark decay modes do not generically evade the constraint. Instead, the additional decay products can increase the radiation density around BBN and strengthen constraints through the primordial helium fraction and \(\Delta N_{\rm eff}\).

This result applies to thermally equilibrated sub-GeV HNLs and does not define a keV sterile-dark-matter region. It provides a lifetime benchmark for the heavy regime, but the board classified the finding as contested because the session did not independently reproduce the underlying calculation.

\subsection{HNL decays through an axion-like particle}

The 2024 study arXiv:2410.06970 considers HNL decays dominated by an ALP and a neutrino. Its abstract reports a new cosmologically consistent region approximately
\[
1\,\mathrm{MeV}\lesssim m_N\lesssim 1\,\mathrm{GeV},
\]
with active--sterile mixing strength approximately
\[
10^{-9}\lesssim \text{mixing strength}\lesssim10^{-6}.
\]
The earlier HNL decay can relax ordinary BBN limits, while the later ALP decay remains constrained by cosmological and astrophysical considerations.

This is a model-dependent MeV--GeV HNL region, not a standard keV sterile-neutrino dark-matter band. The quoted interval was obtained from the paper metadata recorded in the session and was not reproduced numerically.

\subsection{Minimal non-resonant production and additional interactions}

The self-interaction study arXiv:2112.00758 reports that minimal Dodelson--Widrow production supplying all of the dark matter is excluded by X-ray constraints in its benchmark analysis and is also disfavored by structure formation. The same study considers effective non-standard active-neutrino self-interactions. Depending on the interaction strength, the correct-abundance region can shift toward smaller or larger mixing or become wider. A subdominant sterile-neutrino fraction further relaxes the X-ray constraint.

The study reports no structure-formation problem from the additional interactions in its benchmark framework. This statement is conditional on that interaction model and production history, and does not yield a universal \((m_s,\sin^2 2\theta)\) interval. The board marked the finding contested because it was summarized from extracted text rather than checked through a newly performed parameter scan.

\subsection{Resonant production and mixed cold--warm dark matter}

The 2023 study arXiv:2306.16532 examines resonantly produced mixed sterile-neutrino dark matter. Its calculation evolves the sterile-neutrino distribution with a Boltzmann treatment that includes a density potential proportional to the lepton asymmetry. The study assesses viability through the formation of satellite-host subhalos rather than through a complete BBN abundance fit. The session recorded that the required initial lepton number is many orders of magnitude above the baryon-to-photon ratio while remaining below the quoted CMB and BBN detectability limits.

For the \(7.1\,\mathrm{keV}\) benchmark, the study reports models retaining at least 60 Milky Way satellite subhalos. These include nearly all-sterile dark matter at an astrophysical mixing
\[
\sin^2(2\theta_{\rm astro})\simeq 7\times10^{-10},
\]
with viable values approximately below
\[
\sin^2(2\theta_{\rm astro})\lesssim10^{-9}.
\]



Figure 11 (not in record) shows the sampled mixed sterile-neutrino models in terms of astrophysical mixing, sterile fraction, and satellite-subhalo abundance in arXiv:2306.16532. It bears on viability by displaying how the inferred X-ray mixing and the sterile fraction correlate with the structure-formation criterion. The figure is attributed to that paper and is not a new plot generated in the session.



Figure 12 (not in record) shows the corresponding satellite-subhalo viability assessment in arXiv:2306.16532, including the criterion of retaining at least 60 Milky Way satellite subhalos. It bears on the question by identifying which mixed cold--warm models survive the paper's adopted structure-formation requirement. The session did not extract numerical coordinates beyond the reported values above.

The astrophysical mixing is related to the microscopic mixing by
\[
\sin^2(2\theta_{\rm astro})
=
f\,\sin^2(2\theta_{\rm tot}),
\]
where \(f\) is the sterile fraction of the dark matter. Consequently, the quoted \(7\times10^{-10}\) and \(10^{-9}\) values are conditional on \(f\), the production history, and the satellite criterion. They are not universal microscopic mixing limits. The session also recorded that updated scattering rates could lower the allowed maxima. This entire finding was marked contested by the board because the relevant figures and source assumptions were not independently reanalyzed.

\subsection{Scalar-mediated freeze-in and freeze-out}

The mechanism in arXiv:2204.08795 produces keV sterile-neutrino dark matter through two stages. A complex scalar sector is populated through Higgs-portal freeze-in, after which sterile neutrinos are produced by scalar decays and annihilations. The hidden sector can thermalize at a temperature below the Standard Model temperature,
\[
T_\phi<T.
\]
The resulting momentum distribution and free-streaming properties differ from those of oscillation production.



Figure 1 (not in record) is also associated in the session's evidence record with the scalar freeze-in/freeze-out study, but the report does not use a second, independently numbered figure for that paper. The figure information available in the record is insufficient to identify a separate scalar-model plot without risking duplication or misattribution. The scalar mechanism therefore enters the results through its reported two-step production history rather than through an additional figure.

Because the abundance and momentum distribution are changed by the hidden-sector dynamics, this mechanism cannot be represented by the minimal Dodelson--Widrow mass--mixing relation. Its viability must be tested jointly against relic abundance, X-ray decay, and structure-formation constraints. The board classified this finding as contested because no new abundance or free-streaming calculation was performed.

\subsection{Refined resonant production above approximately \(20\,\mathrm{keV}\)}

The 2025 analysis arXiv:2510.01907 revisits resonant production using a refined numerical treatment with dynamically discretized momentum modes. The metadata reports regions consistent with X-ray and Lyman-\(\alpha\) constraints for a lepton-to-entropy ratio
\[
L\gtrsim\mathcal{O}(10^{-3})
\]
and sterile-neutrino mass
\[
m_s\gtrsim20\,\mathrm{keV}.
\]
The study also considers production during lepton-number injection. In that regime, some parameter sets can shorten the free-streaming length and reduce the mixing required to match the observed dark-matter abundance.

These values define a conditional viable region in the refined resonant-production model, not a production-independent band. The session did not extract the associated mixing interval or compare the numerical contour directly with the 7.1-keV benchmark. The board classified the finding as contested.

\subsection{Affleck--Dine resonant production}

The 2024 paper arXiv:2402.11902 studies Shi--Fuller production driven by a large lepton asymmetry generated in an Affleck--Dine scenario with an \(L\)-ball. It revisits X-ray and Lyman-\(\alpha\) constraints and discusses the possible observational consequences of the asymmetry-generating scenario. The available record does not state a numerical surviving mass--mixing interval, and the paper is not a direct BBN analysis.

Thus, the session identified a motivated resonant-production mechanism but did not establish a quantitative viable mass--mixing region for it.

\section{Discussion}

The session's main physical distinction is between keV sterile-neutrino dark matter and heavier HNLs. In the keV regime, the recorded literature emphasizes relic abundance, radiative decay, momentum distributions, free streaming, and small-scale structure. BBN is not established as the dominant constraint for the long-lived keV candidates discussed here. In the MeV--GeV regime, late decay, temporary matter domination, energy injection, and additional radiation can directly modify primordial abundances and \(\Delta N_{\rm eff}\).

The minimal Dodelson--Widrow all-dark-matter scenario is not viable in the cited self-interaction benchmark because of X-ray constraints and is also disfavored by structure formation. This does not exclude sterile-neutrino dark matter in general. Resonant production can alter the momentum distribution through a lepton asymmetry; mixed cold--warm models can reduce the sterile fraction; non-standard self-interactions can change the abundance--mixing relation; and scalar-mediated production can decouple abundance and momentum distribution from the minimal oscillation relation. Each of these possibilities introduces additional model assumptions, so none supplies a universal mass--mixing contour.

The 7.1-keV benchmark and the \(m_s\gtrsim20\,\mathrm{keV}\) refined resonant region should not be directly merged. The former uses a mixed-component model and a satellite-subhalo criterion, while the latter is summarized as satisfying X-ray and Lyman-\(\alpha\) constraints under a specified lepton-asymmetry history. The session did not perform a common calculation with a shared structure-formation baseline, relic-density treatment, or radiative-decay prescription. Their reported viability therefore represents two conditional regions rather than one continuous allowed band.

The BBN results for HNLs likewise cannot be projected onto the keV dark-matter plane. The minimal analysis arXiv:2410.07343 uses \((m_4,\lvert U_{a4}\rvert^2)\) and includes flavor-dependent decay and injection effects, while the ALP and dark-decay studies alter the cosmology by adding new decay channels and radiation components. The ALP scenario opens an approximately \(1\,\mathrm{MeV}\)--\(1\,\mathrm{GeV}\) region with mixing strength approximately \(10^{-9}\)--\(10^{-6}\), whereas the dark-decay study reports that additional radiation can strengthen rather than relax BBN constraints. These statements concern different models and should not be combined into a single HNL exclusion or discovery region.

Several limitations constrain the strength of the conclusions. The autonomous survey failed during full-text extraction. Although direct downloads succeeded for several papers, the session ended before all chunks and figure captions could be read. No numerical BBN contour reproduction, keV parameter scan, free-streaming calculation, lifetime calculation, or joint X-ray--structure analysis was completed. Most board findings were marked contested, and the reported values should therefore be treated as source-reported conditional results rather than independently verified calculations.

\section{Conclusions}

The session established the following limited conclusions. First, the directly relevant minimal BBN analysis arXiv:2410.07343 treats the heavy-neutrino mass and one flavor-specific mixing as its only parameters and includes temporary matter domination plus meson, photon, and neutrino injection. Its qualitative BBN exclusion occupies an intermediate-mixing regime, but the numerical electron-, muon-, and tau-mixing boundaries were not recovered. Second, thermally equilibrated sub-GeV HNLs are reported in arXiv:2507.12270 to require \(\tau_N\lesssim0.02\,\mathrm{s}\), while dark decays can strengthen constraints through extra radiation. Third, an ALP-cascade model in arXiv:2410.06970 reports a conditional region of approximately \(1\,\mathrm{MeV}\)--\(1\,\mathrm{GeV}\) and mixing strength \(10^{-9}\)--\(10^{-6}\).

For keV dark matter, minimal non-resonant all-dark-matter production is excluded by X-rays in the cited self-interaction study and is also disfavored by structure formation. Conditional viable possibilities remain in resonant production, mixed cold--warm models, non-standard self-interactions, and scalar-mediated production. The recorded 7.1-keV mixed model includes nearly all-sterile dark matter near \(\sin^2(2\theta_{\rm astro})\simeq7\times10^{-10}\) and viable values below approximately \(10^{-9}\) under its satellite criterion. The refined resonant analysis reports X-ray- and Lyman-\(\alpha\)-compatible regions for \(L\gtrsim\mathcal{O}(10^{-3})\) and \(m_s\gtrsim20\,\mathrm{keV}\), but no corresponding mixing interval was recovered.

The ranked open directions raised by the session are:

\begin{enumerate}
\item Reproduce the electron-, muon-, and tau-mixing BBN contours of arXiv:2410.07343, including temporary matter domination and decay-injection effects.
\item Construct a conditional keV sterile-neutrino scan combining relic abundance, radiative-decay X-rays, momentum distribution, free streaming, and a specified structure-formation baseline.
\item Compare refined resonant-production histories across the 7.1-keV mixed-component benchmark and the \(m_s\gtrsim20\,\mathrm{keV}\) lepton-asymmetry-driven region.
\item Extend the HNL analysis to ALP-cascade and dark-decay scenarios with explicit lifetime, abundance, decay-temperature, and radiation-density calculations.
\item Quantify the abundance--mixing relation in non-standard neutrino self-interaction and scalar freeze-in/freeze-out models.
\end{enumerate}

\section*{References}

\textit{BBN Constraint on Heavy Neutrino Production and Decay} --- Yu-Ming Chen, Yue Zhang.

\textit{No Hiding in the Dark: Cosmological Bounds on Heavy Neutral Leptons with Dark Decay Channels} --- P. S. Bhupal Dev, Quan-feng Wu, Xun-Jie Xu.

\textit{Relaxing Limits from Big Bang Nucleosynthesis on Heavy Neutral Leptons with Axion-like Particles} --- Frank F. Deppisch, Tomás E. Gonzalo, Chayan Majumdar, Zhong Zhang.

\textit{Resonant production of sterile neutrino dark matter with a refined numerical scheme} --- Kentaro Kasai, Masahiro Kawasaki, Kai Murai.

\textit{Affleck-Dine leptogenesis scenario for resonant production of sterile neutrino dark matter} --- Kentaro Kasai, Masahiro Kawasaki, Kai Murai.

\textit{Exploring resonantly produced mixed sterile neutrino dark matter models} --- Emma L. Horner, Francisco Mungia Wulftange, Isabella A. Ianora, Chad T. Kishimoto.

\textit{Sterile neutrino dark matter production in presence of non-standard neutrino self-interactions: an EFT approach} --- Cristina Benso, Werner Rodejohann, Manibrata Sen, Aaroodd Ujjayini Ramachandran.

\textit{Freeze-in and freeze-out of sterile neutrino dark matter} --- Rupert Coy, Michael A. Schmidt.

\textit{keV sterile neutrino Dark Matter} --- Alexander Merle.

\section*{Data and code availability}
All outputs of this session are collected in the project folder \href{file:///Users/nramachandra/hep-genesis-testing/projects/hep-testing-research}{\texttt{/\allowbreak{}Users/\allowbreak{}nramachandra/\allowbreak{}hep-genesis-testing/\allowbreak{}projects/\allowbreak{}hep-testing-research}}.
The session produced 8 data files; the supplement lists them.
Full documentation of how these results were obtained is in the supplementary material (\texttt{supplement.pdf}: implementation details, the tool-call log, the research board with every finding and its evidence, figures not shown here, and the code index) and the verbatim transcript (\texttt{transcript.pdf}), generated alongside this report.


\end{document}
